% REVISED: canonical-target reframe 2026-06-04
% REVISED: hostile-review action plan 2026-06-04
% REVISED: JLEO compression pass 2026-06-06 — full cost-recall grid, frontier
% figure, baseline/case-holdout tables, and operating-point figures migrate to
% Online Supplement E (\ref{supp:bid_cost}); appendix keeps the sequence
% statement, cost denominators, frontier support, and the algorithm.
% C-COMPRESS-2 (story pass 2026-06-06): folded to ~1.5pp. Appendix keeps the
% sequence statement + cost-denominator table (D1--D8) + frontier-support
% paragraph; the gatekeeping algorithm, baselines/timing/case-holdout detail,
% and operating-point logic move to Online Supplement E (\ref{supp:bid_cost}).
\subsection{Sequential Gatekeeping and the Cost-Recall Frontier}
\label{app:forensic_sequence_submission}

The sequential architecture separates two tasks that are often conflated.
The award-layer stage ranks firms using fields available before any bid
microdata are recovered; the bid-layer stage then constructs richer
forensic features only for the firms that survive the first stage. We
state the sequence and measure the recovery footprint each stage opens,
so the comparison reads as an evidence-allocation design rather than a
claim that thin records substitute for full bid files. The output is a
forensic-priority queue, not a liability finding. Throughout, ``cost''
denotes a \emph{data-processing footprint}---the firms, tender-items, bid
rows, and buyer-item-year cells that must be opened to compute the
bid-layer features---never a monetary agency cost, which we do not
observe. Discriminating performance is adjudication-anchored: positives
are the $\valCostNpos$ always-loser cobidder firms that link to a
CADE-adjudicated case, against the complete-Imhof pool of $16{,}772$
firms. The operating points are retrospective cost-footprint designs
conditional on the validated incumbent ranking, not fully prospective
deployment tests.

Three prediction scores are defined over the pool: an \textbf{award
score} $s_i = \log(1+T_i)$, computable from routine administrative fields
with no bid values; a \textbf{bid score}, a random forest on
within-tender bid-distribution moments requiring the full bid record of
each opened tender-item; and a \textbf{combined score} on the union. We
compare seven ranking rules distinguished by which layer they consume and
whether they triage before opening bid microdata---award-only, bid-only,
joint full-observability (a hypothetical upper-bound ceiling), sequential
award\,$\to$\,bid and award\,$\to$\,combined, FL14\,$\to$\,bid, and a
random survivor-pool floor. The full rule definitions, the bid-feature
dictionary, the random-forest specification, and the four validation
designs are in Online Supplement~\ref{supp:bid_cost}.

\subsubsection{Cost Denominators}
\label{app:joint_scoring_submission}

Bid-layer recovery is costlier at scale because within-tender moments
require the \emph{full} set of bids in each opened tender-item, not only
the focal firm's: opening a survivor firm means opening every tender-item
that firm participated in, in full. We therefore measure the recovery
footprint along several denominators rather than a single headline
number. Table~\ref{tab:G_denoms} defines them.

\begin{table}[htbp]\centering
\begin{adjustbox}{max width=\textwidth,center}
\begin{threeparttable}
\caption{Cost denominators for the recovery footprint.}
\label{tab:G_denoms}
\scriptsize
\begin{tabular}{llr}
\toprule
ID & Denominator (full-observability value) & Value \\
\midrule
D1 & Firms opened ($=|\mathcal{S}|$) & $16{,}772$ \\
D2 & Survivor firm-tender-item rows (with multiplicity) & $163{,}264$ \\
D3 & Opened tender-items (de-duplicated; recovery unit) & $116{,}112$ \\
D4 & Opened bid rows, full-tender (processing footprint) & $3{,}393{,}915$ \\
D5 & Survivor bid rows only (lower bound) & $391{,}610$ \\
D6 & Buyer $\times$ item-group $\times$ year cells touched & $58{,}936$ \\
D7 & Buyers (PBU codes) touched & $1{,}244$ \\
D8 & Item-groups (2-digit) touched & $90$ \\
\bottomrule
\end{tabular}
\begin{tablenotes}\scriptsize
\item \textit{Source:} \texttt{full\_observability\_costs.csv};
denominator definitions in \texttt{table\_G\_cost\_denominator\_definitions.csv}.
Values are the full-pool footprint ($K_1=16{,}772$); the grid in Online
Supplement~\ref{supp:bid_cost} reports them at each $K_1$.
\item \textit{D4 is the realistic processing footprint:} within-tender
Imhof moments require the entire tender record, so D4 counts all bids
in every opened tender-item, not only the survivors' own bids (D5).
\item \textit{Not observed (no value computed):} LANCES export units
(D9); legal or subpoena-style access requests (D10); monetary cost in
reais. \textit{Analyst-hours} are reported only as an illustrative
processing-burden proxy (D11), never as a measured agency cost.
\end{tablenotes}
\end{threeparttable}
\end{adjustbox}
\end{table}

The firm count (D1) and the bid-row footprint (D4) diverge sharply, and
the divergence is the central caution. Survivors are high-participation
firms: at the $K_1=2000$ operating point the survivor pool is about
$12\%$ of the firms (an $88\%$ firm reduction), yet it opens about $67\%$
of the bid rows (only a $33\%$ bid-row reduction). The firm-level
reduction therefore overstates the saving; the tender-item and bid-row
denominators are the honest processing-footprint measures.

\subsubsection{Cost-Recall Frontier and Findings}
\label{app:gatekeeping_algorithm_submission}

The cost-recall grid for the sequential award\,$\to$\,bid rule, sliced at
queue depth $k=500$ across the first-stage cut $K_1$, traces a frontier
that tracks close to the joint full-observability upper bound while
opening fewer tender-items but never crosses it. True positives are
non-monotone in $K_1$ ($124$ at $K_1=1{,}000$ versus $116$ at
$K_1=2{,}000$), one sign that no $K_1$ is optimal, and the sequential
rule approaches but does not exceed the joint ceiling ($124/133 = 93\%$
at $K_1=1{,}000$). The sequential rule also sits well above a same-size
random survivor floor, confirming that it is the award ranking, not the
mere act of opening $K_1$ firms, that carries the signal: the
award-ranked survivor pool captures most positives \emph{before} any bid
rerank, so the bid layer \emph{reorders} survivors into a
forensic-priority queue rather than discovering positives the award gate
missed. The marginal contribution of the costly bid features is ordering,
an evidence-allocation gain, not a detection gain. The full grid (all
rules, all $K_1$, $k\in\{100,250,500,1000\}$ with every cost
denominator), the recall-versus-tender-item frontier figure, the
gatekeeping algorithm, the timing-disciplined and case-holdout exercises,
and the operating-point logic are in Online
Supplement~\ref{supp:bid_cost}. Two qualifications carry the
interpretation: a timing-disciplined award gate remains informative
(strict-timing AUC $0.684$ on the training-AL pool against $0.761$
full-window), while the sequential rule's within-tender moments cannot be
re-derived on a pre-period basis and the timing test is therefore
reported as infeasible; and performance is case-fragile---the largest
case supplies $32.0\%$ of positives and $45.4\%$ of the top-$500$ true
positives, so holding it out lowers recall on smaller cases materially.
No single $(K_1,k)$ pair dominates: the right operating point depends on
agency capacity and on whether the priority is recall or precision, so
the frontier is a menu, not a cutoff. The estimated ranking is not a
portable cartel score; the transferable object is the decomposition
framework.

\subsubsection{Limitations}
\label{app:limitations_submission}

Cost is a processing footprint, not money. Firm reduction overstates the
saving (an $88\%$ firm reduction at $K_1=2000$ is only a $33\%$ bid-row
reduction); analyst-hours are illustrative only; labels are limited and
false positives are unlabeled, so precision is a lower bound on the
queue's triage value; performance is case-fragile; the operating point is
capacity-dependent with no dominating choice; and the sequential timing
test is blocked, so only the award-only gate has a timing-disciplined
evaluation.
